\clearpage
\thispagestyle{empty}
\refstepcounter{section}
\label{app:software-worked-example}

\begin{center}
\textbf{APPENDIX B}

\vspace{\baselineskip}

\textbf{Worked Calculation for the Software-Dependency Case}
\end{center}

\vspace{3\baselineskip}

\subsection{Purpose and Selected Scenario}

This appendix traces one reported software-domain result through the frozen
GBBRPM equations. The case uses the real Express 4.18.2 dependency topology,
the frozen OSV advisory snapshot, unit susceptibility on every retained
dependency relation, and the controlled uniform transmission setting
\(\tau=0.50\). This setting is selected because it demonstrates both
intermediate-node recursion and multi-source convergence at Express. It is a
sensitivity setting, not an empirically calibrated software transmission
coefficient.

The original deps.dev relation points from a dependent package to the
dependency it uses. For propagation, the relation is reversed so that a local
disturbance in a dependency can travel toward packages that depend on it. The
calculation below therefore uses dependency-to-dependent direction. Package
names are shortened for readability, but every value corresponds to the exact
package version reported in Chapter IV.

\subsection{Local Disturbance and Relation Inputs}

The seven evidence-backed package versions have the local disturbances shown
in Table~\ref{tab:app-software-inputs}. Each relation uses \(S_{ij}=1\) and
the selected uniform sensitivity value \(\tau_{ij}=0.50\).

\begin{table}[H]
\caption{Software Inputs Used in the Worked Calculation}
\label{tab:app-software-inputs}
\centering
\small
\begin{tabularx}{\textwidth}{@{}L c c L@{}}
\toprule
Package version & Severity & \(B_i\) & Relevant outgoing relation(s) \\
\midrule
\texttt{qs@6.11.0} & Moderate & 0.50 & \texttt{body-parser}, Express \\
\texttt{send@0.18.0} & Low & 0.25 & \texttt{serve-static}, Express \\
\texttt{body-parser@1.20.1} & High & 0.75 & Express \\
\texttt{path-to-regexp@0.1.7} & High & 0.75 & Express \\
\texttt{cookie@0.5.0} & Low & 0.25 & Express \\
\texttt{serve-static@1.15.0} & Low & 0.25 & Express \\
\texttt{express@4.18.2} & Moderate & 0.50 & Root dependent package \\
\bottomrule
\end{tabularx}
\end{table}

Packages without advisory-backed local disturbance have \(B_i=0\). Because
their upstream states are also zero in this frozen case, their contributions
to the nodes calculated below are zero and their \((1-Q_{ij})\) terms equal
one. They can therefore be omitted from the displayed products without
changing the result.

\subsection{Step 1: Initialize the Disturbed Dependency Nodes}

The selected upstream disturbed packages have no positive incoming
contribution in this scenario. Their accumulated states therefore equal their
local disturbances:

\[
R_{\texttt{qs}}=0.50,
\qquad
R_{\texttt{send}}=0.25,
\qquad
R_{\texttt{path-to-regexp}}=0.75,
\qquad
R_{\texttt{cookie}}=0.25.
\eqtag{B-1}
\]

These are comparative normalized disturbances derived from the declared
advisory-severity mapping; they are not exploitation probabilities.

\subsection{Step 2: Evaluate the Intermediate Dependent Packages}

The contribution from \texttt{qs} to \texttt{body-parser} is

\[
\begin{aligned}
Q_{\texttt{qs},\texttt{body-parser}}
 &=S\tau R_{\texttt{qs}} \\
 &=(1.00)(0.50)(0.50) \\
 &=0.25.
\end{aligned}
\eqtag{B-2}
\]

Combining this incoming contribution with the local disturbance
\(B_{\texttt{body-parser}}=0.75\) gives

\[
\begin{aligned}
R_{\texttt{body-parser}}
 &=1-(1-0.75)(1-0.25) \\
 &=1-(0.25)(0.75) \\
 &=0.8125.
\end{aligned}
\eqtag{B-3}
\]

Similarly, the contribution from \texttt{send} to \texttt{serve-static} is

\[
Q_{\texttt{send},\texttt{serve-static}}
 =(1.00)(0.50)(0.25)=0.125.
\eqtag{B-4}
\]

The receiving package also has local disturbance
\(B_{\texttt{serve-static}}=0.25\), so

\[
\begin{aligned}
R_{\texttt{serve-static}}
 &=1-(1-0.25)(1-0.125) \\
 &=1-(0.75)(0.875) \\
 &=0.34375.
\end{aligned}
\eqtag{B-5}
\]

Equations B-3 and B-5 demonstrate current-state recursion: the states later
transmitted to Express contain both local and inherited effects.

\subsection{Step 3: Calculate the Contributions Entering Express}

At \(\tau=0.50\) and \(S=1\), the six nonzero contributions entering
Express are shown in Table~\ref{tab:app-software-contributions}.

\begin{table}[H]
\caption{Nonzero Contributions Entering Express at \(\tau=0.50\)}
\label{tab:app-software-contributions}
\centering
\small
\begin{tabular*}{0.86\textwidth}{@{\extracolsep{\fill}}lrr}
\toprule
Dependency source \(i\) & Current \(R_i\) & \(Q_{i,\mathrm{Express}}=R_i(1)(0.50)\) \\
\midrule
\texttt{body-parser} & 0.812500 & 0.406250 \\
\texttt{cookie} & 0.250000 & 0.125000 \\
\texttt{path-to-regexp} & 0.750000 & 0.375000 \\
\texttt{qs} & 0.500000 & 0.250000 \\
\texttt{send} & 0.250000 & 0.125000 \\
\texttt{serve-static} & 0.343750 & 0.171875 \\
\bottomrule
\end{tabular*}
\end{table}

For example, the first contribution uses the accumulated intermediate state:

\[
Q_{\texttt{body-parser},\mathrm{Express}}
 =(1.00)(0.50)(0.8125)=0.40625.
\eqtag{B-6}
\]

Using only the original \(B_{\texttt{body-parser}}=0.75\) would instead give
0.375 and would incorrectly discard the contribution already inherited from
\texttt{qs}.

\subsection{Step 4: Aggregate Local and Inherited State at Express}

Express has its own advisory-derived local disturbance \(B=0.50\). Applying
the multiplicative-complement rule to the six nonzero incoming contributions
gives

\[
\begin{aligned}
R_{\mathrm{Express}}
={}&1-(1-0.50)
\times(1-0.40625)(1-0.125)(1-0.375)\\
&\quad\times(1-0.25)(1-0.125)(1-0.171875)\\
={}&1-(0.50)(0.59375)(0.875)(0.625)\\
&\quad\times(0.75)(0.875)(0.828125)\\
={}&1-(0.50)(0.1764643192)\\
={}&0.9117678404\approx0.9118.
\end{aligned}
\eqtag{B-7}
\]

This independently reproduces the Chapter IV result for Express at
\(\tau=0.50\). At this setting Express becomes the highest-ranked node, while
all states remain within \([0,1]\). The example establishes traceability of
the software instantiation and the effect of the controlled transmission
sweep; it does not validate prediction of exploitation or operational
incidents.
